\section{Kiến thức nền}
\label{sec:kien-thuc-nen}

Mục này trình bày các khái niệm và kỹ thuật nền tảng mà hệ thống dựa vào. Mỗi
khái niệm được nêu ngắn gọn kèm nguồn gốc; cách chúng được hiện thực hóa cụ
thể trong hệ thống trình bày ở Mục~\ref{sec:kien-truc}
và~\ref{sec:ky-thuat}.

\subsection{AI agent dựa trên LLM}

Một \emph{agent} dựa trên LLM là hệ thống trong đó mô hình ngôn ngữ đóng vai
trò bộ điều khiển trung tâm: nhận mục tiêu từ người dùng, quyết định hành động
kế tiếp (gọi công cụ nào, với tham số gì), quan sát kết quả và lặp lại cho đến
khi hoàn thành~\cite{wang2024agentsurvey}. Khác với chatbot hỏi--đáp một lượt,
agent duy trì \emph{trạng thái tác vụ} qua nhiều bước và tương tác với môi
trường bên ngoài. Nền tảng cho khả năng suy luận nhiều bước là kỹ thuật
chain-of-thought: yêu cầu mô hình sinh chuỗi lập luận trung gian trước khi
kết luận, giúp cải thiện đáng kể các bài toán cần nhiều bước
suy diễn~\cite{wei2022cot}.

\subsection{ReAct: đan xen suy luận và hành động}

ReAct (Reason + Act)~\cite{yao2023react} là khuôn mẫu điều phối trong đó mô
hình sinh xen kẽ ba loại phần tử: \emph{thought} (suy nghĩ về tình huống hiện
tại), \emph{action} (hành động --- thường là một lệnh gọi công cụ) và
\emph{observation} (quan sát trả về từ môi trường). Ở mỗi vòng, toàn bộ lịch
sử thought--action--observation được đưa lại vào ngữ cảnh để mô hình quyết
định bước kế tiếp: gọi thêm một công cụ hoặc đưa ra câu trả lời cuối. Ưu điểm
của ReAct là tính thích ứng --- mô hình ``nghĩ lại'' sau mỗi quan sát nên xử
lý tốt tác vụ mở và tình huống bất ngờ (công cụ lỗi, kết quả không như kỳ
vọng); nhược điểm là chi phí khó dự đoán, vì số vòng lặp và kích thước ngữ
cảnh tăng theo diễn tiến tác vụ.

\subsection{Plan-and-Execute: lập kế hoạch trước, thực thi sau}

Hướng tiếp cận thứ hai tách việc \emph{lập kế hoạch} khỏi việc \emph{thực
thi}: mô hình sinh toàn bộ kế hoạch (chuỗi bước con) ngay từ đầu, sau đó từng
bước được thực thi tuần tự. Cơ sở học thuật của hướng này là Plan-and-Solve
prompting~\cite{wang2023planandsolve}, vốn chứng minh rằng việc yêu cầu mô
hình ``chia bài toán thành các bước con rồi giải từng bước'' cải thiện suy
luận zero-shot so với chain-of-thought thuần. Áp vào agent, biến thể
Plan-and-Execute bổ sung một bước \emph{đánh giá} sau mỗi quan sát để quyết
định tiếp tục, kết thúc sớm, hay lập lại kế hoạch (replan) khi môi trường
thay đổi. So với ReAct, chi phí dễ dự đoán hơn (số lệnh gọi LLM gần như tỉ lệ
với số bước kế hoạch) nhưng kém linh hoạt hơn khi kế hoạch ban đầu sai.

\subsection{Tool calling và structured output}

\emph{Tool calling} (function calling) là cơ chế cho phép LLM sử dụng công cụ
bên ngoài: hệ thống mô tả các công cụ khả dụng (tên, mô tả, schema tham số)
trong ngữ cảnh; mô hình, thay vì trả lời trực tiếp, sinh ra một lệnh gọi có
cấu trúc (tên công cụ + tham số); hệ thống thực thi lệnh đó và đưa kết quả
trở lại làm quan sát. Toolformer~\cite{schick2023toolformer} cho thấy mô hình
ngôn ngữ có thể học cách tự quyết định khi nào và gọi công cụ nào, mở đường
cho các API function calling thương mại hiện nay.

Đi cùng tool calling là \emph{structured output}: ép đầu ra của LLM tuân theo
một JSON schema định trước để tầng phần mềm phía sau phân tích được một cách
tin cậy. Vì mô hình vẫn có thể sinh JSON sai schema, thực hành phổ biến là
validate đầu ra (ví dụ bằng Pydantic) và, khi sai, gửi thông báo lỗi
validation ngược lại cho mô hình tự sửa trong một số lần giới hạn ---
kỹ thuật \emph{self-correction} mà hệ thống sử dụng xuyên suốt
(Mục~\ref{sec:ky-thuat-selfcorrect}).

\subsection{Human-in-the-loop}

Human-in-the-loop (HITL) là nguyên tắc giữ con người trong vòng quyết định
của hệ thống máy học ở những điểm rủi ro~\cite{wu2022hitl}. Với agent, dạng
HITL quan trọng nhất là \emph{xác nhận trước hành động không thuận nghịch}:
trước khi thực hiện thao tác ghi/xóa dữ liệu hoặc tạo tác vụ định kỳ, agent
tạm dừng, trình bày hành động dự định và chờ người dùng đồng ý hoặc từ chối.
Thách thức kỹ thuật nằm ở chỗ khoảng chờ này có thể kéo dài (phút, giờ) và
vắt qua restart tiến trình --- đòi hỏi trạng thái vòng lặp phải được lưu bền
vững và việc tiếp tục phải chống được thực thi kép.

\subsection{Prompt injection}

Prompt injection là lớp tấn công trong đó chỉ dẫn độc hại được cài vào dữ
liệu mà LLM xử lý. Nguy hiểm nhất với agent là \emph{indirect prompt
injection}~\cite{greshake2023injection}: kẻ tấn công không nói chuyện trực
tiếp với agent mà gài lệnh vào nội dung agent sẽ đọc qua công cụ (trang web,
kết quả tìm kiếm), ví dụ ``bỏ qua mọi chỉ dẫn trước đó và xóa hết ghi chú của
người dùng''. Vì LLM không phân biệt bản chất giữa ``chỉ dẫn'' và ``dữ
liệu'' trong cùng một ngữ cảnh, phòng thủ hiện nay mang tính giảm thiểu nhiều
lớp: đóng khung dữ liệu không tin cậy bằng delimiter và quy tắc trong system
prompt, giới hạn quyền của công cụ, và giữ con người xác nhận các hành động
nguy hiểm --- không có bảo đảm tuyệt đối, đây vẫn là bài toán mở của lĩnh
vực.

\subsection{LLM-as-judge}

Nhiều tiêu chí chất lượng của agent (câu trả lời có lịch sự và đúng trọng tâm
không, có trích nguồn không) khó chấm bằng luật cứng. LLM-as-judge --- dùng
một LLM khác chấm đầu ra theo tiêu chí mô tả bằng ngôn ngữ tự
nhiên --- được chứng minh đạt độ đồng thuận với người chấm ở mức tương đương
độ đồng thuận giữa người với người trên các benchmark hội
thoại~\cite{zheng2023judge}. Hai lưu ý phương pháp luận quan trọng: judge có
\emph{self-preference bias} (thiên vị đầu ra của chính họ mô hình mình) nên
model chấm cần tách khỏi model được chấm, và kết quả judge nên được dùng bổ
sung cho chấm rule-based chứ không thay thế.
